% Template:     Informe LaTeX
% Documento:    Archivo principal
% Versión:      8.4.8 (12/07/2026)
% Codificación: UTF-8
%
% Autor: Pablo Pizarro R.
%        pablo@ppizarror.com
%
% Manual template: [https://template-latex.github.io/informe]
% Licencia MIT:    [https://opensource.org/licenses/MIT]

% CREACIÓN DEL DOCUMENTO
\documentclass[
	spanish, % Idioma: spanish, english, etc.
	letterpaper, oneside
]{article}

% INFORMACIÓN DEL DOCUMENTO
\def\documenttitle {Apunte Física Computacional II}
\def\documentsubtitle {}
\def\documentsubject {}

\def\documentauthor {Nombre del autor}
\def\coursename {}
\def\coursecode {}

\def\universityname {Universidad de Concepción}
\def\universityfaculty {Facultad de Ciencias Físicas y Matemáticas}
\def\universitydepartment {Departamento de Física}
\def\universitydepartmentimage {departamentos/fcfm}
\def\universitydepartmentimagecfg {height=1.57cm}
\def\universitylocation{Concepción}

% INTEGRANTES, PROFESORES Y FECHAS
\def\authortable {
	\begin{tabular}{ll}
		Integrantes:
		& \begin{tabular}[t]{l}
			Joaquín Parra \\
			Ignacio Falcón \\
      Alvaro Osses \\
      Benjamin Fuentes
		\end{tabular} %\\
 % 	Profesor:
 % 	& \begin{tabular}[t]{l}
 % 		Profesor 1
 % 	\end{tabular} \\
 % 	Auxiliar:
 % 	& \begin{tabular}[t]{l}
 % 		Auxiliar 1
 % 	\end{tabular} \\
 % 	Ayudantes:
 % 	& \begin{tabular}[t]{l}
 % 		Ayudante 1 \\
 % 		Ayudante 2
 % 	\end{tabular} \\
 % 	\multicolumn{2}{l}{Ayudante de laboratorio: Ayudante 1} \\
 % 	& \\
 % 	\multicolumn{2}{l}{Fecha de realización: \today} \\
 % 	\multicolumn{2}{l}{Fecha de entrega: \today} \\
 % 	\multicolumn{2}{l}{\universitylocation}
	\end{tabular}
}

% IMPORTACIÓN DEL TEMPLATE
\input{template}

% INICIO DE PÁGINAS
\begin{document}

% PORTADA
\templatePortrait

% CONFIGURACIÓN DE PÁGINA Y ENCABEZADOS
\templatePagecfg

% RESUMEN O ABSTRACT
\begin{abstractd}
  El siguiente documento presenta un apunte para el curso de Física Computacional II, basado en las notas de clase y presentaciones del ramo dictado por Dr. Roberto Navarro durante el segundo semestre de los años 2024 y 2025, y las notas de clase redactadas por Pedro Contreras. 
\end{abstractd}

% TABLA DE CONTENIDOS - ÍNDICE
\templateIndex

% CONFIGURACIONES FINALES
\templateFinalcfg

% ======================= INICIO DEL DOCUMENTO =======================

%\input{example} % Ejemplo, se puede borrar
\section{Cálculo Numérico}

\subsection{Errores Computacionales}

Si usted ha intentado realizar una suma de dos números decimales mediante algún lenguaje de programación, podrá haber notado que el resultado no es exactamente correcto
\begin{sourcecode}{python}{Suma de dos decimales en Python}
  a = 0.1
  b = 0.2
  c = 0.3

  print(f"{a + b == c}")
  print(f"{a + b}")
\end{sourcecode}
Cuyo output será

\begin{sourcecode}{plaintext}{Output}
  False
  0.30000000000000004
\end{sourcecode}

Este error se debe a la manera en que las computadoras deben representar números. En general, los humanos representamos los números en base 10, por ejemplo, escribimos el 325 como
\begin{equation*}
  325 = 3\cdot 10^{2} + 2\cdot 10^{1} + 5\cdot 10^{0} = [325]_{10}
\end{equation*}
Pero los computadores utilizan bits (0 y 1), por lo que estos representan números en base 2, es decir:
\begin{equation*}
  325 = 2^{8} + 2^{6} + 2^{2} + 2^{0} = [101000101]_{2}
\end{equation*}

Ahora, para representar número reales en el sistema binario, se deben incluir exponentes negativos a la suma de exponentes. En la representación binaria, los números se separan por un exponente y una parte fraccionaria (mantisa) las cuales son separadas por un punto ($.$). Por ejemplo,

\begin{align*}
  [2.145]_{10} &= 2\cdot 10^{0} + \frac{1}{8} \\
               &= 2^{1} + 2^{-3} \\
               &= [10.001]_{2} \\
               &= [1.0001]_{2} \cdot 2^{1}
\end{align*}

Este ejemplo en particular funciona bien porque la mantisa del númer decimal resulta ser un exponente de 2, veamos que ocurre para un número más complejo:

\begin{align*}
  [0.1]_{10} &= [1]_{2}/[1010]_{2} \\
             &= [0.0\overline{0011}]_{2} \\
             &= [0.00011001100110011 \cdots] \\
             &= 2^{-4} + 2^{-5} + 2^{-8} + 2^{-9} + 2^{-12} + 2^{-13} + \cdots \\
             &= 2^{-4}\cdot (1+ 2^{-1} + 2^{-4} + 2^{-5} + \cdots)
\end{align*}
Obtenemos una serie infinita de potencias de dos. Como en la practica los computadores tienen memoria finita es necesario truncar la serie. En general, la mantisa recibe 52 bits, por lo que la serie se trunca hasta el termino $2^{-52}$. en efecto:
\begin{align*}
  [0.1]_{10} &= 2^{-4}\cdot(1+ 2^{-1} + 2^{-4} + 2^{-5} + \cdots + 2^{-52})
\end{align*}

Este número no será \textit{exactamente} igual a $0.1$, si no que será igual a $0.1 \cdot (1+ \epsilon)$. Este será el caso para la gran mayoría de los calculos hechos computacionalmente, el resultado final tendrá una pequeña divergencia con respecto al valor real, es decir 
$$\bar{a} = a\cdot (1+ \epsilon) $$
donde $\bar{a}$ y $a$ son los valores calculado y real, respectivamente. La diferencia entre el número deseado y el resultante se denomina \textbf{error computacional}. Dado que las computadores no pueden físicamente lidiar con cantidades infinitesimalmente pequeñas, muchos de los métodos de cálculo adaptados a la programación que veremos a lo largo del apunte serán basados en aproximaciones, por esta razón es importante aprender a identificar fuentes de errores numericos y sus ordenes de magnitud.

\subsection{Derivadas numéricas}

Sabemos del calculo diferencial que la derivada de una función arbitraria $f$ evaluada en un punto arbitrario $x_{0}$  es dada por
\begin{equation}
  f^{\prime}(x_{0}) = \lim_{h \to 0} \frac{f(x_{0}+h) - f(x_{0})}{h}
\end{equation}
pero ya sabemos que necesitamos que $h$ sea un número finito, por lo que estamos obligados a seguir anaizando esta definición. De la serie de taylor de $f$ evaluada en un punto $a + h$, $h > 0$, tenemos que:

\begin{equation}
  f(a + h) = f(a) + f^{\prime}(a)h + \frac{1}{2!}f^{\prime \prime}(a)h^{2} + \frac{1}{3!}f^{\prime \prime \prime}(a)h^{3} + \cdots
\end{equation}
Podemos reemplazar los terminos de orden superior por un $f(\xi)h^{2}$, $a \leq \xi \leq a + h$ (Teorema del valor medio), luego
\begin{equation}
  f(a + h) = f(a) + f^{\prime}(a)h + \frac{1}{2!}f^{\prime \prime}(\xi)h^{2}
\end{equation}
\begin{equation} \label{eq:derad}
  f^{\prime}(a) = \frac{f(a + h) - f(a)}{h} - \frac{1}{2!}f^{\prime \prime}(\xi)h
\end{equation}

Aquí, hemos obtenido una expresión finíta para la derivada de una función evaluada en un punto, denominada \textit{derivada adelantada}. Ahora, si buscamos el error de la ecuación \eqref{eq:derad} vemos que
\begin{equation}
  Error =  \left| f^{\prime}(a) - \frac{f(a+h) - f(a)}{h} \right| = \frac{1}{2!}\left| f^{\prime \prime}(\xi) \right| h
\end{equation}
en notación $\mathcal{O}$ de Landau, 
\begin{equation}
  Error = \frac{1}{2!}\left| f^{\prime \prime}(\xi) \right| h = \mathcal{O}(h)
\end{equation}
Esto quiere decir que el error de truncamiento de la derivada adelantada será del orden de $h$, como se puede observar en la figura \ref{fig:adnterr}, el error de la derivada decrece linealmente con el orden de magnitud de $h$ hasta aproximadamente $h \sim 10^{-8}$, pasado ese punto $h$ se vuelve demasiado pequeño para la computadora y predomina el error computacional. Por esta razón, a la hora de implementar este, o cualquiera de los métodos expuestos en este apunte, es conveniente utilizar un $h \sim 10^{-5}$. 

\insertimage[\label{fig:adnterr}]{./img/derivadas/adnterr.pdf}{scale = 0.75}{Error de la derivada adelantada de $f(x) = \sin(x)$ en $x = 0.3$ para distintos valores de $h$}

\subsubsection{Otros métodos de derivación numérica}

Si bien la derivada adelantada es útil debido a lo simple de su implementación, muchas veces queremos utilizar métodos cuyo error sea de un mayor orden de $h$, es decir, menor.

Similarmente a como obtuvimos la expresión \eqref{eq:derad}, podemos obtener una derivada \textbf{retrasada}, en efecto, consideramos la serie de taylor de $f(a-h)$,
\begin{align}
  f(a - h) &= f(a) - f^{\prime}(a)h + \frac{1}{2!}f^{\prime \prime}(\xi)h^{2} \\
  f^{\prime}(a) &= \frac{f(a) - f(a-h)}{h} + \frac{1}{2!}f^{\prime \prime}(\xi)h \label{eq:deriret}
\end{align}

Si consideramos las series de taylor de $f(a+h)$ y $f(a-h)$ expandidas haste el tercer orden,
\begin{align}
  f(a+h) &= f(a) + f^{\prime}(a)h + \frac{1}{2!}f^{\prime \prime}(a)h^{2} + \frac{1}{3!}f^{\prime \prime \prime}(\xi)h^{3} \\
  f(a-h) &= f(a) - f^{\prime}(a)h + \frac{1}{2!}f^{\prime \prime}(a)h^{2} - \frac{1}{3!}f^{\prime \prime \prime}(\xi)h^{3}
\end{align}
Ahora si restamos obtenemos,
\begin{align}
  f(a+h) - f(a-h) &= 2f^{\prime}(a)h + \frac{2}{3!}f^{\prime \prime \prime}(\xi)h^{3} \\
  2f^{\prime}(a)h &= f(a+h) - f(a-h) - \frac{2}{3!}f^{\prime \prime \prime}(\xi)h^{3} \\
  f^{\prime}(a) &= \frac{f(a+h) - f(a-h)}{2h} - \frac{1}{6}f^{\prime \prime \prime}(\xi)h^{3} \label{eq:dercen}
\end{align}

Así, obtenemos una expresión para la \textbf{derivada centrada}, con error del orden de $h^{3}$.

Podemos obtener muchos otros métodos de derivación numérica por medios similares, en general el proceso consiste en jugar con series de taylor hasta poder despejar el termino $f^{\prime}(a)$. 

\subsubsection{Derivadas para datos discretos}

En general, uno no suele tratar con funciones definidas analíticamente, si no que tendremos conjuntos de datos relacionados a alguna magnitud física, si quisieramos eventualmente encontrar la derivada de la "función" que definen estos datos, debemos aprender a tratar con puntos no necesariamente equidistantes.

Similarmente a caso anterior, dado un conjunto de puntos $x_{k}$ y una función continua y diferenciable tal que $f(x_{k}) = f_{k}$, definimos la distancia entre dos puntos como $h_{k} = x_{k} - x_{k-1}$. Podemos obtener una primera aproximación simplemente tomando la pendiente entre dos puntos, esto es
\begin{equation} \label{eq:firstdisc}
  \frac{df_{k}}{dx_{k}} \approx \frac{ f_{k}-f_{k-1}}{x_{k}-x_{k-1}}
\end{equation}
Esto es analogo a la derivada atrasada encontrada en la sección anterior, y se puede demostrar que tendrá un error del orden de $h_{k}$. Notar además que si imponemos que los nodos estén igualmente espaciados, es decir, $x_{k}-x_{k-1} = h = cte.$, la ecuación \eqref{eq:firstdisc} recupera la ecuación \eqref{eq:deriret}.

Ahora, estamos incentivados a encontrar un método de derivación discreta con error de menor orden, para ello consideramos la serie de taylor de $f_{k+1}$ y $f_{k-1}$ 
\begin{align}
  f_{k+1} &= f_{k} + \fracderivat{f_{k}}{x}h_{k+1} + \frac{1}{2!}\fracnderivat{f_{k}}{x}{2}h_{k+1}^{2} + \frac{1}{3!}\fracnderivat{f(\xi_{1})}{x}{3}h_{k+1}^{3} \label{eq:kplus} \\
    f_{k-1} &= f_{k} - \fracderivat{f_{k}}{x}h_{k} + \frac{1}{2!}\fracnderivat{f_{k}}{x}{2}h_{k}^{2} - \frac{1}{3!}\fracnderivat{f(\xi_{2})}{x}{3}h_{k}^{3} \label{eq:kmin}
\end{align}

con $x_{k-1} \leq \xi_{2} \leq x_{k} \leq \xi_{1} \leq x_{k+1}$. Despejando $\fracderivat{f_{k}}{x}$ en ambas expresiones obtenemos,
\begin{align}
  \fracderivat{f_{k}}{x}h_{k+1} &= f_{k+1} - f_{k} - \frac{1}{2!}\fracnderivat{f_{k}}{x}{2}h_{k+1}^{2} - \fracnderivat{f(\xi_{1})}{x}{3}h_{k+1}^{3} \label{eq:kone} \\
    \fracderivat{f_{k}}{x}h_{k} &= f_{k} - f_{k-1} + \frac{1}{2!}\fracnderivat{f_{k}}{x}{2}h_{k}^{2} + \fracnderivat{f(\xi_{2})}{x}{3}h_{k}^{3} \label{eq:ktwo}
\end{align}
Multiplicando \eqref{eq:kone} y \eqref{eq:ktwo} por $h_{k}/h_{k+1}$ y $h_{k+1}/h_{k}$, respectivamente, 
\begin{align}
  \fracderivat{f_{k}}{x}{}h_{k} &= \frac{(f_{k+1} - f_{k})h_{k}}{h_{k+1}} - \frac{1}{2!}\fracnderivat{f_{k}}{x}{2}h_{k}h_{k+1} - \fracnderivat{f(\xi_{1})}{x}{3}h_{k+1}^{2}h_{k} \\
    \fracderivat{f_{k}}{x}{}h_{k+1} &= \frac{(f_{k} - f_{k-1})h_{k+1}}{h_{k}} + \frac{1}{2!}\fracnderivat{f_{k}}{x}{2}h_{k}h_{k+1} + \fracnderivat{f(\xi_{2})}{x}{3}h_{k}^{2}h_{k+1}
\end{align}
Sumando,
\begin{equation}
  [h_{k+1} + h_{k}]\fracderivat{f_{k}}{x} = \frac{(f_{k+1} - f_{k})h_{k}}{h_{k+1}} + \frac{(f_{k} - f_{k-1})h_{k+1}}{h_{k}} + \frac{1}{3!} \left[ f^{\prime \prime \prime}(\xi_{2})h_{k} - f^{\prime \prime \prime}(\xi_{1})h_{k+1} \right]h_{k}h_{k+1}
\end{equation}
Finalmente,

\begin{equation}
  \fracderivat{f_{k}}{x} = \frac{(f_{k+1} - f_{k})}{h_{k+1} + h_{k}}\frac{h_{k}}{h_{k+1}} + \frac{(f_{k} - f_{k-1})}{h_{k+1} + h_{k}}\frac{h_{k+1}}{h_{k}} + \frac{1}{3!} \left[ \frac{ f^{\prime \prime \prime}(\xi_{2})h_{k} - f^{\prime \prime \prime}(\xi_{1})h_{k+1}}{h_{k+1} + h_{k}} \right]h_{k}h_{k+1}
\end{equation}
Ahora, por el teorema del valor medio, podemos reemplazar el último termino entre corchetes por una función evaluada en un punto $\xi$, con $\xi_{2} \leq \xi \leq \xi_{1}$, en efecto,

\begin{align}
  \fracderivat{f_{k}}{x} &= \frac{(f_{k+1} - f_{k})}{h_{k+1} + h_{k}}\frac{h_{k}}{h_{k+1}} + \frac{(f_{k} - f_{k-1})}{h_{k+1} + h_{k}}\frac{h_{k+1}}{h_{k}} + \frac{1}{3!} \fracnderivat{f_{k}(\xi)}{x}{3} h_{k}h_{k+1} \\
  &= \frac{(f_{k+1} - f_{k})}{h_{k+1} + h_{k}}\frac{h_{k}}{h_{k+1}} + \frac{(f_{k} - f_{k-1})}{h_{k+1} + h_{k}}\frac{h_{k+1}}{h_{k}} + \mathcal{O}(h_{k}h_{k+1})
\end{align}

Así, logramos encontrar una expresión para la derivada de una función definida punto a punto con nodos no necesariamente equidistantes, con error de orden cuadratico. Puede comprobar que si imponemos que los nodos sean equidistantes ($h_{k} = h_{k+1} = h$), se recupera la expresión \eqref{eq:dercen} para la derivada centrada (\textbf{Demuestrelo!}). Más adelante estudiaremos métodos más avanzados para tratar con distribuciones de datos con interpolaciones y splines. Para terminar, cabe mencionar que la mayor desventaja de este método es que, al utilizar 3 puntos, aplicandola a una serie de puntos, la expresión numérica de la derivada tendrá dos puntos menos que la serie original. Por esta razón, en la practica se implementa una derivada adelantada para los puntos iniciales de un set de datos, y una atrasada para los finales; la derivada centrada se utiliza solo para los puntos intermedios.


\section{Breve introducción a Python}

Python es un lenguaje de alto nivel con una sintaxis enfocada en la legibilidad y la simpleza. Los principios fundamentales del lenguaje pueden ser encontrados \textit{dentro} del mismo, pruebe ejecutar el código \ref{thispy}

\begin{sourcecode}[\label{thispy}]{python}{Easter egg}
  import this
\end{sourcecode}

\textbf{Spoiler}: Si lo hace, el output será el \textbf{Zen de Python}
\begin{sourcecode}[\label{thisout}]{plaintext}{Zen de Python}
Beautiful is better than ugly.
Explicit is better than implicit.
Simple is better than complex.
Complex is better than complicated.
Flat is better than nested.
Sparse is better than dense.
Readability counts.
Special cases aren't special enough to break the rules.
Although practicality beats purity.
Errors should never pass silently.
Unless explicitly silenced.
In the face of ambiguity, refuse the temptation to guess.
There should be one-- and preferably only one --obvious way to do it.
Although that way may not be obvious at first unless you're Dutch.
Now is better than never.
Although never is often better than *right* now.
If the implementation is hard to explain, it's a bad idea.
If the implementation is easy to explain, it may be a good idea.
Namespaces are one honking great idea -- let's do more of those!
\end{sourcecode}




\newpage
\section{Introducción a la Búsqueda de Ceros}

Usualmente, cuando trabajamos con ecuaciones, una de nuestras metas es encontrar el valor de cierta variable o también llamada incógnita, para esto, desarrollamos una serie de pasos o ''algoritmo'' que seguir para encontrar la respuesta. Un buen ejemplo de esto es considerar el caso de una ecuación cuadrática:

\begin{equation}
ax^2 + bx^2 +c = 0
\end{equation}

Es bien sabido que, dado los coeficientes a, b y c, existe una fórmula general que seguir, la conocida como ecuación cuadrática o \textit{Ecuación de Bashkara}:

\begin{equation}
x = \frac{-b \pm \sqrt{b^2-4ac}}{2a}
\end{equation}

Al introducir los respectivos valores de los coeficientes, podemos encontras las dos soluciones a este problema.

Sin embargo, no siempre se puede encontrar una ecuación para los valores de incógnitas en toda ecuación, un ejemplo a esto es el siguiente:

\begin{equation}
\sin x + x = 0
\end{equation}

A pesar de lo simple de nuestro problema (puesto que es la suma entre una función trigonométrica y la incógnita), ya no es posible encontrar una ecuación de forma analítica que exprese valores exactos de la soluciones para ''x'', así, la pregunta es ¿Qué podemos hacer en estos casos?

Para esto, se desarrollan los \textbf{Métodos para hallar ceros de una función}

Estos, se basan en una idea principal, convertir la ecuación en la forma:

\begin{equation}
    f(x) = 0
\end{equation}

Esto, con la finalidad de que ahora nuestro trabajo, será encontrar el valor para el cual, ''x'' cumple con ser un ''cero'' de la función.

Para encontrar este cero, es importante establecer un ''intervalo'' en donde estemos trabajando (puesto que la computadora no es capaz de analizar, en caso de trabajar en $\mathbb{R}$, la recta $]-\infty,+ \infty[$, ya que es infinita), para esto, tenemos que considerar un intervalo en donde deba existir una solución.

Sin embargo, ¿Cómo podemos asegurar una solución?

Para ello, existe el \textbf{Teorema de Bolzano}, el cual establece que:

\begin{quote}
\textit{Para una función continua en un intervalo cerrado $[a,b]$, tal que $f(a) \cdot f(b) < 0$, entonces debe existir $c \in [a,b]$ tal que $f(c) = 0$}
\end{quote}

Así, ya podemos concluir que el método con el que trabajemos esté siendo usado en un intervalo donde la solución exista.

Es importante considerar que debido a que utilizaremos métodos numéricos, la solución que encontraremos nunca será del todo correcta, sino, tan solo una aproximación a la respuesta real o analítica.

Por ello, es importante tener en cuenta \textit{¿Qué tan cerca queremos o podemos estar de nuestra respuesta real?}. Para esto, es importante tener en cuenta el uso de una medida de ''Tolerancia''.

La tolerancia, se puede definir como el intervalo más pequeño sobre el cual vamos a realizar el algoritmo de búsqueda de nuestro cero. Cuando nuestros intervalos sean suficientemente pequeños, le indicaremos al programa que detenga el algoritmo y se quede con el valor o respuesta dentro del intervalo encontrado.

Así, si la tolerancia (denotada como $\delta$) cumple que: $f(x_i) < \delta$, entonces la solución numérica encontrada será $x_i$.

Esto conlleva a la cuestión ¿Cómo podemos medir el error que se encuentra dentro de nuestro método?


% Arreglar esto!!!
Para esto, es importante analizar \textit{qué tan rápido aumenta el error}, tema del cual se verá dentro de cada método.

A continuación, exploraremos algunos de los métodos para encontrar ceros.


\subsection{Método de la Secante}

Consideremos la siguiente ecuación:

\begin{equation}\label{ec_biseccion_cos}
\cos x - 2x - x^2 = 0
\end{equation}

Podemos notar que esta ecuación (debido a los términos que posee) a simple vista no tiene una solución analítica que calcular, sin embargo, podemos utilizar métodos numéricos.

El método de la bisección se basa en analizar entre que valores estamos trabajando o buscando la solución. Consideremos para este caso, el intervalo $[0,\frac{\pi}{2}]$.

Es posible notar, evaluando dentro de los límites de los intervalos, que:

\begin{itemize}
\item $f(0) = \cos 0 - 0 -0 > 0$
\item $f(\frac{\pi}{2}) = \cos \frac{\pi}{2} - \pi - \left(\frac{\pi}{2}\right)^2 < 0$
\end{itemize}

Así, podemos notar que existe una región dentro del intervalo, donde la función que establecimos cambia de signo, bien, que \textbf{El producto entre la función evaluada en los límites del intervalo es menor a cero}.

A partir de esto, hay que recordar el \textit{Teorema de Bolzano}, el cual afirma que:

\begin{quote}
Si en una función continua, existe un cambio de signo en los límites de un intervalo cerrado, entonces esta posee al menos una raíz.
\end{quote}

Así, podemos confirmar que existe una solución a la ecuación (\ref{ec_biseccion_cos}) dentro del intervalo $[0,\frac{\pi}{2}]$, a partir de esto, consideraremos realizar el siguiente proceso:

Dividiremos el intervalo actual en dos intervalos con la misma distancia, sean $[0,\frac{\pi}{4}]$, $[\frac{\pi}{4},\frac{\pi}{2}]$. Y evaluaremos los límites sobre el primero de estos. De aquí, tendremos 3 posibles respuestas:

\begin{itemize}
    \item Si $f(0) \cdot f(\frac{\pi}{4}) < 0$: La solución debe encontrarse en este intervalo, por lo que nos centraremos en trabajar aquí.
    \item Si $f(0) \cdot f(\frac{\pi}{4}) = 0$: La solución está en uno de los límites de los intervalos, como en el proceso anterior el resultado fue negativo, entonces podemos confirmar que la solución es $x= \frac{\pi}{4}$
    \item Si $f(0) \cdot f(\frac{\pi}{4}) > 0$: La solución no está dentro de este intervalo, así, debemos trabajar con el intervalo siguiente ($\frac{\pi}{4}, \frac{\pi}{2}$).
\end{itemize}

Una vez realizado este análisis, encontraremos la solución o bien, el intervalo donde se encuentra. A partir de esto, podemos seguir realizando el mismo algoritmo para encontrar intervalos más pequeños en donde nuestra solución se debe encontrar, esto, hasta alcanzar el límite o tolerancia establecida.

La tolerancia, dentro de este capítulo, se define como el intervalo más pequeño sobre el cual vamos a realizar el algoritmo, cuando nuestros intervalos sean suficientemente pequeños, le indicaremos al programa que detenga el algoritmo y se quede con el valor intermedio del intervalo encontrado.

De esta forma, podemos catalogar el método de la bisección de forma general

Considerar una ecuación de la forma $f(x) = 0$ en el intervalo $[a,b]$

\begin{itemize}
    \item Analizamos el producto de la función $f(x)$ evaluada en los límites ($x = a,b$), esto es $f(a) \cdot f(b)$
    \item Si el producto es negativo, el valor "x" se encuentra dentro de este intervalo, por lo que podemos continuar con el procedimiento.
    \item Dividimos el intervalo en dos partes, sea $[a,c], [c,b]$
    \item Analizamos el producto de la función evaluada en los límites del primer intervalo
    \item Si este producto es negativo, la solución debe encontrarse en este intervalo, por que dividimos y volvemos a aplicar el mismo análisis
    \item Repetimos este proceso hasta que la distancia del intervalo sea menor que nuestro valor de tolerancia.

\end{itemize}

Podemos observar un ejemplo de este a través del siguiente código:
\begin{lstlisting}
f = lambda x: np.cos(x)
tolerancia = 1e-10
while abs(a-b) > tolerancia:
    if f(a)*f(b) < 0:
        c = (b+a)/2
        if f(a)*f(c) < 0:
            b = c
        else: 
            a = c
print(c)
\end{lstlisting}

Para nuestro caso, hemos de elegir la función \textit{cos(x)}, si consideramos el intervalo $[0,2]$, podemos encontrar el cero que se posiciona en $x = \frac{\pi}{2}$, donde bajo este código, obtenemos un valor de $x \approx 1.5707963$, bastante cercano al buscado.

\subsection{Método de Newton}

Consideraremos buscar el cero de nuestra función f(x) a través del uso de la serie de Taylor.

Para esto, tomemos en nuestro algoritmo la función $f(x)$, si desarrollamos en segundo orden la serie de Taylor para el cero, se puede notar que:

\begin{equation}
f(x_0)= f(x) + (x - x_0) \cdot f'(x) + (x - x_0)^2 \cdot \frac{f''(x)}{2} + ...
\end{equation}

Pero por definición, la función evaluada en nuestro punto $x_0$ debe ser cero, así, si aproximamos la serie:

\begin{equation}
0 = f(x) + (x-x_0) \cdot f'(x)
\end{equation}

Podemos despejar así $x_0$, donde:

\begin{equation}
x_0 = x - \frac{f(x)}{f'(x)}
\end{equation}

Esto, nos ofrece un algoritmo que seguir, a modo de poder aproximarse a través de iteraciones hasta superar la tolerancia establecida.

Consideraremos un punto inicial, tal como el valor inicial del intervalo dentro del cual estamos buscando el cero, para este caso, será $x_n$, así, la iteración del algoritmo nos dará un valor cercano al cero, siendo este denotado $x_{n+1}$, así:

\begin{equation}
x_{n+1} = x_n - \frac{f(x_n)}{f'(x_n)}
\end{equation}

Un ejemplo de la implementación de este algoritmo es el siguiente, consideramos una función $f(x) = x^3 -  5x$, y su derivada $f'(x) = 3x^2 - 5$, así:

\begin{lstlisting}
tolerancia = 1e-12

f = lambda x: x**3 -5*x # Se define la función a utilizar
fprima = lambda x: 3*x**2 - 5

def met_newton(f,a,b,tolerancia):
    zeros = []
    int = np.linspace(a,b,10)
    print(int)
    for i in range(int.size-1): 

        if f(int[i])*f(int[i+1]) <= 0:
            a,b = int[i:i+2]
            fa, fprima_a = f(a), fprima(a)
            m = 0
            while (abs(fa/fprima_a )) >= tolerancia:
                b = a - fa/fprima_a
                a = b
                fa, fprima_a = f(a), fprima(a)
                m +=1
            print("Existe un cero",b,"con iteraciones de",m,"veces")


met_newton(f,-3,3,tolerancia)

\end{lstlisting}

Es importante mencionar que el método de Newton es uno de los métodos más rápidos a la hora de encontrar ceros, esto, gracias al uso de la derivada de la función.

A su vez, notar que si la derivada se acerca a cero, la cantidad de iteraciones comenzará a aumentar, puesto que la velocidad con la que avanza el algoritmo comienza a disminuir.


\subsection{Método de la Secante}

\textit{Nota del autor: Este capítulo lo hice en el portafolio de física compu 2}

¿Qué ocurre si no podemos calcular u obtener la derivada de la función?

Esto es común en problemas donde la derivada puede no existir en ciertas partes, por lo que es importante analizar casos de este estilo.

Para esto, nace el método de la secante, el cual se basa en la idea de aproximar la derivada de la función, de tal forma que, si consideramos diferencias lo suficientemente pequeñas entre $x_n$ y $x_{n-1}$, entonces:

\begin{equation}
    f'(x_n) \approx \frac{f(x_n) - f(x_{n-1}}{x_n - x_{n-1}}
\end{equation}

Por lo que el algoritmo a iterar constaría de la siguiente estructura:

\begin{equation} \label{metodo_secante}
x_{n+1} = x_n - \frac{x_n - x_{n-1}}{f(x_n) - f(x_{n-1})}f(x_n)
\end{equation}

Un ejemplo al utilizar este método es el siguiente:

\begin{lstlisting}

tolerancia = 1e-5

f = lambda x: x**3 -5*x # Se define la función a utilizar

    # Se define el método de la secante
def met_secante(f,a,b,tolerancia):
    zeros = []
    int = np.linspace(a,b,10)
    print(int)
    for i in range(int.size-1): 

        if f(int[i])*f(int[i+1]) <= 0:
            a,b = int[i:i+2]
            fa,fb = f(a), f(b)
            m = 0
            while (abs(fb*(b-a)/(fb-fa)) >=tolerancia):
                bold=b 
                b = b - fb*(b-a)/(fb-fa) 
                fa=fb 
                a=bold
                fb=f(b) 
                m +=1
            print("Existe un cero",b,"con iteraciones de",m,"veces")
            zeros.append(b)
    return zeros

met_secante(f,-3,3,tolerancia)

\end{lstlisting}

Ahora bien, es importante considerar \textbf{¿Cuál es la proporción de error dentro del método de la secante?}




% FIN DEL DOCUMENTO
\end{document}
